\chapter{Rétro-ingénierie pour la compilation avec un SDK}
\label{chap:t6}
% Entreprise : Beemo Technologie — Volet « maîtriser la chaîne de production » (2/2)

\section{Contexte et dette technique initiale}
% TODO : un programme doit être compilé avec un \gls{sdk} donné, mais la
% compilation échoue / n'est pas documentée. C'est un cas typique de dette
% (build cassé, savoir perdu, dépendance opaque). Enjeu business : sans
% build, pas de livraison.
Cette tâche est sensiblement plus \textbf{courte} que celles menées chez Free
Mobile, mais elle n'en est pas moins \textbf{déterminante} dans la démarche de
réduction de la dette technique. Elle porte sur l'un des logiciels de Beemo,
\textbf{l'agent VSX}, et vise à reprendre le contrôle sur une partie de sa
\textbf{chaîne de production}. Sa portée dépasse le simple correctif : il s'agit
d'\textbf{automatiser} et d'\textbf{harmoniser} la méthode de fabrication de cet
agent — ce qui la rend, à mes yeux, aussi structurante que la réalisation du
framework de test étalée sur près d'un an et demi.

Avant cette mission, mettre à jour l'agent VSX supposait une manipulation
manuelle et fragile : on \textbf{poussait la base de code vers une machine
distante}, sur laquelle elle était compilée avec le \textbf{DSDK} — le
\gls{sdk} fourni par Broadcom pour la réalisation de \emph{daemons} (processus
d'arrière-plan) —, avant de pouvoir être testée.

La dette technique est apparue avec le \textbf{passage à la version~8 du DSDK},
nécessaire pour supporter la version~8 de l'agent. Dans sa configuration de
base, le DSDK impose le flag de compilation C++ \texttt{-Wall -Werror}, qui
\textbf{convertit l'ensemble des avertissements en erreurs} bloquantes. Or
l'agent est écrit en C++11, avec une \textbf{absence quasi totale de bonnes
pratiques} de génie logiciel — certains fichiers dépassent les
\textbf{50~000~lignes}. Le volume d'avertissements produit est dès lors
\textbf{colossal} et intraitable en l'état : il met un frein considérable à la
capacité de l'équipe à \textbf{livrer des mises à jour en continu}.

Une rustine existait bien — des modifications appliquées directement sur la
machine de compilation pour neutraliser ce comportement — mais elle n'était
\textbf{ni reproductible ni documentée}. Ce savoir reposait sur le développeur
\textbf{que je remplace} : son départ menaçait de l'emporter avec lui.

\section{Objectif de la mission}
% TODO : rendre la compilation fonctionnelle avec le \gls{sdk}, ce qui a
% nécessité une démarche de \gls{retro-ingenierie}.
La mission a poursuivi deux objectifs complémentaires :
\begin{itemize}
  \item \textbf{s'affranchir de la machine de compilation distante} en rendant
        le \emph{build} reproductible localement, à l'aide d'un conteneur
        \textbf{Docker} embarquant la \gls{toolchain} du DSDK ;
  \item \textbf{rendre les correctifs du SDK reproductibles et documentés}, en
        substituant aux modifications manuelles des \textbf{fichiers de patch}
        versionnés, appliqués automatiquement lors de la construction de
        l'image.
\end{itemize}
Atteindre ces deux objectifs supposait au préalable de \textbf{comprendre le
fonctionnement interne du DSDK} — une démarche de \gls{retro-ingenierie}, le
SDK n'exposant aucun moyen standard d'agir sur ses flags de compilation.

\section{Démarche de rétro-ingénierie}
% TODO : c'est le cœur de cette tâche — détailler la méthode :
% - analyse des artefacts (binaires, symboles, \gls{abi}/\gls{api}, headers)
% - reconstitution de la \gls{toolchain} / des dépendances manquantes
% - outils utilisés (objdump, nm, strings, ldd, débogueur... à adapter)
% - itérations d'essais-erreurs
Le cœur de la tâche a consisté à \textbf{retrouver où et comment le DSDK injecte
ses flags de compilation}. Cette mécanique n'est \textbf{pas exposée} par le
SDK : il n'existe aucune option officielle pour assouplir le couple
\texttt{-Wall -Werror}. Il a donc fallu \textbf{analyser le système de build du
DSDK lui-même} pour localiser le point précis où ces flags sont posés, puis
déterminer comment les neutraliser sans dénaturer le reste de la
\gls{toolchain}.

La difficulté principale n'a pas été d'identifier la modification à apporter,
mais de \textbf{trouver le bon point d'injection} : le DSDK \textbf{n'est pas
conçu pour être modifié}, et un patch mal placé compromet l'intégrité de la
chaîne. Une fois ce point localisé, la suite s'est révélée nettement plus
simple : il a suffi de cibler les fichiers concernés et d'en dériver des
correctifs reproductibles.

\section{Réalisation}
% TODO : extrait de configuration de build / commande / patch ayant
% débloqué la compilation (commenté).
La solution repose sur un \textbf{Dockerfile} qui reconstitue l'environnement de
compilation auparavant confiné à la machine distante, puis applique les
correctifs du SDK de façon déterministe. Chaque modification est isolée dans un
fichier \texttt{.patch} versionné et appliquée via la commande \texttt{patch}
lors de la construction de l'image — là où elle était jusqu'alors saisie à la
main, sans trace.

\begin{lstlisting}[language=bash,caption={Reconstitution de la toolchain et application reproductible des correctifs du DSDK dans le Dockerfile (extrait illustratif).},label={lst:t4}]
# Image de compilation embarquant la toolchain Broadcom (DSDK v8),
# auparavant disponible uniquement sur une machine distante.
FROM broadcom/dsdk:8

# Application reproductible des correctifs du SDK : chaque .patch est
# versionne et documente, en lieu et place des modifications manuelles
# faites jusqu'ici directement sur la machine de compilation.
COPY patches/ /tmp/patches/
RUN patch -p1 -d "${DSDK_ROOT}" < /tmp/patches/0001-assouplir-werror.patch

# Le build de l'agent VSX devient ainsi local et rejouable a l'identique.
WORKDIR /src
RUN make agent-vsx
\end{lstlisting}

\section{Difficultés rencontrées}
% TODO : absence de doc, incompatibilités de versions, opacité du SDK.
La principale difficulté a tenu à l'\textbf{opacité du DSDK} et à
l'identification du point d'injection des patchs, évoquée plus haut. Une seconde
est apparue en aval : l'un des correctifs appliqués au SDK a eu un
\textbf{effet de bord sur la certification} de l'un des \emph{packages} générés,
qui n'est plus accepté en l'état. C'est la raison pour laquelle la tâche
demeure \textbf{encore en cours} : le \emph{build} local et reproductible est
acquis, mais la \textbf{validation du package} reste à fiabiliser.

\section{Résultats et bénéfice business}
% TODO : build reproductible et documenté = dette résorbée, déblocage du
% produit, capitalisation du savoir pour l'équipe.
La mission a transformé une \textbf{procédure manuelle, distante et non
documentée} en un \emph{build} \textbf{local, reproductible et versionné}. Les
correctifs du SDK, jadis tributaires d'un seul développeur, sont désormais
\textbf{capitalisés sous forme de patchs} : leur logique survit à tout départ,
ce qui répond directement au risque de \textbf{perte de connaissance} à l'origine
de la dette. En supprimant le frein imposé par \texttt{-Wall -Werror}, l'équipe
retrouve sa capacité à \textbf{livrer des mises à jour de l'agent VSX en
continu}.

Le chantier restant — résoudre l'impact sur la certification — s'inscrit dans
une perspective plus large : \textbf{automatiser le processus de validation d'un
package par Broadcom}, aujourd'hui \textbf{très chronophage}. À terme, c'est
l'ensemble de la chaîne de production de l'agent qui gagnerait en
\textbf{harmonisation} et en \textbf{autonomie}, prolongeant la résorption de la
dette bien au-delà du seul déblocage de la compilation.
